\begin{frame}{Anexă: funcții codificate în biți}
  Pentru $f:\N\to\N$, scriem:
  \[
    \underbrace{1\cdots1}_{f(0)}0\;
    \underbrace{1\cdots1}_{f(1)}0\;
    \underbrace{1\cdots1}_{f(2)}0\;\cdots
  \]
  Exemplu: $f=2,0,3,\ldots$ produce $110\,0\,1110\cdots$.
  \medskip
  Separatorul 0 ne permite să recuperăm toate valorile.
  \bigskip
  Și invers, orice șir de biți este o funcție $\N\to\N$ care ia doar valorile 0 și 1.
\end{frame}
\note{Prima aplicație este o injecție N\textasciicircum{}N-\textgreater{}2\textasciicircum{}N. Nu este surjectivă: șirurile codificate au infinit de multe zerouri. Injecția în sens invers este includerea funcțiilor cu valori 0 și 1 în funcțiile cu valori naturale. Teorema Cantor--Schröder--Bernstein transformă două injecții în existența unei bijecții. Nu este necesar să construim explicit acea bijecție pentru argumentul de necalculabilitate.}

\begin{frame}{Anexă: biți și numere reale}
  Pentru biții $b_0,b_1,b_2,\ldots$, folosim un real cu cifre ternare 0 sau 2:
  \[
    x=\sum_{i=0}^{\infty}\frac{2b_i}{3^{i+1}}.
  \]
  Șiruri diferite dau numere diferite.
  \bigskip
  În sens invers, codificăm un real prin semn, partea întreagă și o scriere zecimală canonică.
  \bigskip
  \textbf{Două injecții dau o bijecție}\par
  prin teorema Cantor--Schröder--Bernstein.
\end{frame}
\note{Dacă două șiruri diferă prima dată la i, diferența termenului principal este 2/3\textasciicircum{}\{i+1\}, iar coada poate compensa cel mult 1/3\textasciicircum{}\{i+1\}, deci diferența rămâne pozitivă. Pentru real-\textgreater{}N\textasciicircum{}N, codificăm semnul lui x, partea întreagă a lui \textbar{}x\textbar{} și cifrele fracționare, alegând expansiunea fără coadă infinită de 9 și un singur cod pentru zero. Aceasta este o injecție. Împreună cu injecția N\textasciicircum{}N-\textgreater{}2\textasciicircum{}N-\textgreater{}R și Cantor--Schröder--Bernstein obținem \textbar{}N\textasciicircum{}N\textbar{}=\textbar{}R\textbar{}. Folosim teorema numai pentru cardinalitate, nu pretindem o bijecție calculabilă pe toți realii arbitrari.}

